% !PRINTED 5
\weberpage{5}
Die Aufgabe, die uns nun zunächst beschäftigen wird, ist eine doppelte: Es soll durch Einführung neuer Veränderlicher das Differential
\weq{5}{\frac{dx}{\sqrt{f(x)}} \quad \text{oder} \quad \frac{y\,dx-x\,dy}{\sqrt{f(x,y)}}}
in ein anderes von derselben Form transformiert werden, worin aber die Funktion unter dem Quadratwurzelzeichen möglichst vereinfacht und besonders von einer möglichst kleinen Zahl von Parametern abhängig gemacht wird (§ 3, 4, 5). Es soll zweitens das allgemeine Differential (3) oder (4) in andere ähnliche zerlegt werden, in welchen die Funktion $\Phi(x)$ oder $\Phi(x,y)$ möglichst einfach ist (§ 12).

Die Funktion $f(x)$ läßt sich in lineare Faktoren zerlegen (Bd. I, § 33). Wir setzen demnach auch
\[
f(x) = a_0(x-x_1)(x-x_2)(x-x_3)(x-x_4).
\]
Denken wir uns die Größen $x, x_1, x_2, \ldots$ als Abszissen von Punkten auf eine gerade Linie von einem beliebigen Nullpunkte aus aufgetragen, die positiven nach der einen, die negativen nach der anderen Seite (z. B. positiv nach rechts, negativ nach links), so stellen die Differenzen $x-x_1, x-x_2, x-x_3, x-x_4$ die Entfernungen des Punktes $x$ von $x_1, x_2, x_3, x_4$ dar, und zwar positiv gerechnet, wenn $x$ nach der positiven Seite von $x_1, x_2, \ldots$ liegt. In demselben Sinne ist $x_2-x_1$ die Entfernung des Punktes $x_2$ von $x_1$ usf.

Wir bezeichnen die Differenzen auch kürzer durch $(xx_1), (xx_2), \ldots (x_2x_1), \ldots$ Wir benutzen diese Darstellungsweise aber nur zur leichteren Übersicht und schließen auch imaginäre Punkte der Geraden nicht aus.

§ 2. \gesperrt{Doppelverh\"altnisse.}

Das einfachste Mittel zur Vereinfachung des elliptischen Differentials stützt sich auf die lineare Transformation der Funktion $f(x)$ (Bd. I, § 67). Es mögen $\alpha, \beta, \gamma, \delta$ vier beliebige Konstanten bedeuten, deren Determinante
\weq{1}{\alpha\delta - \beta\gamma = r}
von Null verschieden ist. Wir setzen
\weq{2}{x = \frac{\alpha x' + \beta}{\gamma x' + \delta},}
%CONT>>
